% !TEX program = lualatex
% =============================================================================
%  ПОВНИЙ ШАБЛОН ЛИСТА — dstu4163 (ДСТУ 4163:2020)
%  Задіяно всі можливості класу. Щоб вимкнути реквізит, очистіть його
%  значення ({}) або аргумент команди — він зникне разом із відступами.
%  Порожня версія з тими самими примітками — lyst-porozhnii.tex.
% =============================================================================

% 14pt — кегль (12pt | 13pt | 14pt); onehalf — інтервал тексту 1,5 (п. 7.3).
\documentclass[14pt,onehalf]{dstu4163}

% TikZ потрібен лише для умовної емблеми нижче; власний логотип вставляйте
% через \includegraphics[height=17mm]{файл}.
\usepackage{tikz}

\dstusetup{
  % --- Вид бланка та автор --------------------------------------------------
  blank      = corner,        % corner | long | none
  sender     = organisation,  % organisation | person
  align      = center,        % center | flush
  form       = false,         % true — порожні дата/№/«На №» стають лініями

  % --- Реквізити бланка (01–08) ---------------------------------------------
  emblem     = {\tikz{\draw[line width=0.8pt] (0,0) circle (7.5mm);
                      \node[font=\bfseries\small] at (0,0) {ЕДК};}},   % (02) ≤ 17 мм
  parent     = {Умовна асоціація консультантів},          % (03)
  name       = {Товариство з обмеженою відповідальністю
                «Економічні дослідження та консалтинг»},                % (04)
  short-name = {(ТОВ «ЕДК»)},                                          % (04)
  unit       = {Відділ прикладних досліджень},                         % (05)
  address    = {вул. Велика Васильківська, 72,\\ м. Київ, 03150},     % (06)
  phone      = {(044) 123-45-67},
  email      = {office@edk.example},
  web        = {www.edk.example},
  extra      = {IBAN UA00 0000 0000 0000 0000 0000 0},
  code       = {12345678},                                             % (08)

  % --- Реквізити документа ---------------------------------------------------
  date       = {18.11.2026},  % (10) або \dstutoday
  index      = {01-12/348},   % (11)
  ref-index  = {112/04-15},   % (12) лист, на який відповідаємо
  ref-date   = {05.11.2026},  % (12)
  access     = {Для службового користування\\ Прим. № 1},  % (14) зазвичай порожнє
  % doctype  = {Гарантійний лист},  % (09) у звичайних листах не зазначається

  % --- Оформлення -------------------------------------------------------------
  heading-bold = true,        % напівжирний заголовок (п. 7.2)
  page-numbers = true,        % номер з другої сторінки (п. 7.10)
}

% --- (15) Адресати: до чотирьох; частини — через порожній рядок ------------
\addressee{%
  Науково-дослідний центр «Зразок»

  Завідувачу відділу\\ макроекономічного аналізу\\
  Олені \surname{Коваль}}
\addressee{%
  Директору ТОВ «Аналітичні системи»\\
  Василю \surname{Федорчуку}}

\begin{document}

% --- (19) Заголовок до тексту ----------------------------------------------
\heading{Про надання даних щодо ринку оренди житла та умови їх використання}

% --- Звертання до кількох адресатів ----------------------------------------
\salutation{Шановні колеги!}

% --- (20) Текст -------------------------------------------------------------
На Ваш запит від 05.11.2026 № 112/04-15 повідомляємо, що ТОВ «ЕДК» готове
надати знеособлені дані про ставки оренди житла у містах обласного значення
за 2023–2025 роки в розрізі адміністративних районів і типів житла.

Дані сформовано за результатами щоквартального моніторингу оголошень
про оренду. Перед передаванням їх буде агреговано до рівня, що унеможливлює
ідентифікацію окремих об'єктів нерухомості або орендодавців.

Умови використання даних:

1.~Дані використовуються виключно в наукових і аналітичних цілях.

2.~У публікаціях обов'язкове посилання на ТОВ «ЕДК» як на джерело даних.

3.~Передавання даних третім особам без письмової згоди ТОВ «ЕДК» не
допускається.

Просимо до 01.12.2026 підписати угоду про використання даних (додається)
і повідомити контактну особу, відповідальну за їх отримання. Після
надходження підписаної угоди дані буде надіслано протягом п'яти робочих днів
у форматі електронної таблиці на зазначену Вами адресу електронної пошти.

Додатково пропонуємо провести спільний онлайн-семінар для обговорення
методології збирання й очищення даних, а також можливих напрямів
подальшої співпраці.

% --- (21) Додатки: перелік із нумерацією ----------------------------------
\begin{attachments}
  \item Угода про використання даних на 4 арк. у 2 прим.
  \item Опис набору даних і методології на 6 арк. в 1 прим.
\end{attachments}

% --- Заключна формула (необов'язкова; у листах організацій часто без неї) --
\closing{З повагою}

% --- (22) Підписи: один під одним за підпорядкованістю; seal — «М. П.» ----
\signature{Директор}{Ігор \surname{Шевчук}}
\signature[seal]{Начальник відділу\\ прикладних досліджень}{Олег \surname{Марченко}}

% --- (27) Виконавці: до 2–3, внизу останньої сторінки ----------------------
\executor{Мельник Ольга 044 123 45 68, o.melnyk@edk.example}
\executor{Гордієнко Сергій 044 123 45 70}

\end{document}
